\documentclass[11pt,a4paper]{article}
\usepackage[utf8]{inputenc}
\usepackage[T1]{fontenc}
\usepackage[ngerman]{babel}
\usepackage{helvet}
\renewcommand{\familydefault}{\sfdefault}
\usepackage[left=2cm, right=2cm, top=2cm, bottom=2cm]{geometry}
\usepackage{amsmath,amssymb}
\usepackage{array}
\usepackage{booktabs}
\usepackage{tikz}
\usepackage{pgfplots}
\pgfplotsset{compat=1.18}
\usepackage{fancyhdr}

\pagestyle{fancy}
\fancyhf{}
\fancyhead[L]{\small Physik Klasse 10E}
\fancyhead[R]{\small Probeklausur Kinematik -- Variante C}
\fancyfoot[C]{\small Seite \thepage}
\renewcommand{\headrulewidth}{0.4pt}

\setlength{\parindent}{0pt}
\setlength{\fboxsep}{8pt}

% Aufgaben-Befehl im Layout B Stil
\newcounter{aufgabe}
\newcommand{\aufgabe}[2]{%
\stepcounter{aufgabe}%
\vspace{0.5cm}%
\noindent\fbox{\parbox{\dimexpr\textwidth-2\fboxsep-2\fboxrule}{%
\textbf{Aufgabe \theaufgabe: #1}\hfill\textbf{#2 BE}%
}}%
\vspace{0.3cm}%
}

\begin{document}

% === KOPFBEREICH ===
\noindent\fbox{\parbox{\dimexpr\textwidth-2\fboxsep-2\fboxrule}{
\centering
{\Large\bfseries Probeklausur: Kinematik -- Variante C}\\[0.2cm]
{\small Klasse 10E \quad|\quad 90 Minuten \quad|\quad 60 BE}
}}

\vspace{0.4cm}
\begin{tabular}{@{}p{6cm}p{8cm}@{}}
Name: & \hrulefill \\[0.3cm]
Datum: & \hrulefill \\
\end{tabular}

\vspace{0.3cm}

\begin{center}
\small
\begin{tabular}{|c|c|c|c|c|c|c|c|c|c|}
\hline
\textbf{Aufgabe} & 1 & 2 & 3 & 4 & 5 & 6 & 7 & Zusatz & \textbf{Summe} \\
\hline
\textbf{max. BE} & 6 & 8 & 8 & 8 & 8 & 10 & 10 & (2) & 60 \\
\hline
\textbf{erreicht} & & & & & & & & & \\
\hline
\end{tabular}
\end{center}

\vspace{0.3cm}

% === AUFGABEN ===

\aufgabe{Grundlagen}{6}

\begin{enumerate}
\item[a)] Rechne um: 36 km/h = \underline{\hspace{2cm}} m/s \hfill (1 BE)

\item[b)] Ein Ball fällt vom Balkon (Höhe 45 m). Berechne die Fallzeit bis zum Aufprall. \hfill (2 BE)

\vspace{2.5cm}

\item[c)] Ergänze die Tabelle: \hfill (3 BE)

\begin{center}
\begin{tabular}{|l|c|c|}
\hline
\textbf{Diagramm} & \textbf{Steigung bedeutet} & \textbf{Fläche bedeutet} \\
\hline
$s(t)$-Diagramm & & --- \\
\hline
$v(t)$-Diagramm & & \\
\hline
\end{tabular}
\end{center}
\end{enumerate}

\aufgabe{Motorrad-Bremsung}{8}

Ein Motorradfahrer fährt mit 72 km/h durch eine Baustelle. Plötzlich taucht ein Hindernis auf. Der Fahrer bremst mit $a = -4\,\mathrm{m/s^2}$.

\begin{enumerate}
\item[a)] Rechne die Geschwindigkeit in m/s um. \hfill (1 BE)

\vspace{1.2cm}

\item[b)] Berechne die Bremszeit bis zum Stillstand. \hfill (2 BE)

\vspace{2cm}

\item[c)] Berechne den Bremsweg. \hfill (2 BE)

\vspace{2.5cm}

\item[d)] Das Hindernis ist 40 m entfernt. Kann der Fahrer rechtzeitig anhalten? Begründe und berechne den Abstand zum Hindernis. \hfill (3 BE)

\vspace{2.5cm}
\end{enumerate}

\newpage

\aufgabe{Volleyball-Aufschlag}{8}

Ein Spieler wirft den Ball zum Aufschlag senkrecht nach oben mit $v_0 = 10\,\mathrm{m/s}$.

\begin{enumerate}
\item[a)] Berechne die Steigzeit (Zeit bis zum höchsten Punkt). \hfill (2 BE)

\vspace{2.5cm}

\item[b)] Berechne die maximale Steighöhe über der Abwurfposition. \hfill (3 BE)

\vspace{3cm}

\item[c)] Mit welcher Geschwindigkeit kommt der Ball wieder auf Abwurfhöhe an? Begründe ohne Rechnung. \hfill (3 BE)

\vspace{2.5cm}
\end{enumerate}

\aufgabe{Autofahrt in der Stadt}{8}

Das $v(t)$-Diagramm zeigt die Fahrt eines Autos:

\begin{center}
\begin{tikzpicture}
\begin{axis}[
    width=11cm, height=5cm,
    xlabel={$t$ in s},
    ylabel={$v$ in m/s},
    xmin=0, xmax=20,
    ymin=0, ymax=15,
    grid=major,
    axis lines=left,
    xtick={0,4,8,12,16,20},
    ytick={0,3,6,9,12,15}
]
% Schattierte Flächen zur Visualisierung
\addplot[fill=black!10, draw=none] coordinates {(0,0) (4,0) (4,12) (0,0)} \closedcycle;
\addplot[fill=black!10, draw=none] coordinates {(4,0) (16,0) (16,12) (4,12)} \closedcycle;
\addplot[fill=black!10, draw=none] coordinates {(16,0) (20,0) (16,12)} \closedcycle;
% Graph
\addplot[thick, black] coordinates {(0,0) (4,12) (16,12) (20,0)};
% Phasen-Beschriftungen
\node at (axis cs:2,4) {\small A};
\node at (axis cs:10,6) {\small B};
\node at (axis cs:18,4) {\small C};
\end{axis}
\end{tikzpicture}
\end{center}

\begin{enumerate}
\item[a)] Beschreibe die drei Phasen A, B und C in Worten. \hfill (3 BE)

\vspace{2.5cm}

\item[b)] Berechne die Beschleunigung in Phase A. \hfill (2 BE)

\vspace{2cm}

\item[c)] Berechne den Gesamtweg der 20-sekündigen Fahrt. \hfill (3 BE)

\vspace{3cm}
\end{enumerate}

\newpage

\aufgabe{Fallexperimente}{8}

\begin{enumerate}
\item[a)] In einem Vakuumrohr lässt man eine Feder und eine Münze gleichzeitig fallen. Beide kommen gleichzeitig unten an. Erkläre physikalisch, warum das so ist. \hfill (3 BE)

\vspace{3.5cm}

\item[b)] Derselbe Versuch wird ohne Vakuum (in Luft) wiederholt. Erkläre, warum die Münze jetzt zuerst ankommt. Nenne zwei relevante physikalische Größen. \hfill (3 BE)

\vspace{3.5cm}

\item[c)] Ein Regentropfen fällt aus einer Wolke in 2000 m Höhe. Berechne die Geschwindigkeit, die er ohne Luftwiderstand beim Aufprall hätte. Erkläre, warum die tatsächliche Geschwindigkeit viel geringer ist. \hfill (2 BE)

\vspace{3cm}
\end{enumerate}

\aufgabe{Weg-Zeit-Gesetz aufstellen}{10}

Ein Zug fährt mit konstanter Geschwindigkeit $v = 25\,\mathrm{m/s}$ durch einen Bahnhof. Zum Zeitpunkt $t = 0$ befindet er sich an der Position $s_0 = 50\,\mathrm{m}$ (gemessen vom Bahnhofseingang).

\begin{enumerate}
\item[a)] Stelle das Weg-Zeit-Gesetz $s(t)$ für den Zug auf. \hfill (2 BE)

\vspace{2cm}

\item[b)] Berechne, wo sich der Zug nach 8 Sekunden befindet. \hfill (2 BE)

\vspace{2cm}

\item[c)] Wann passiert der Zug die Position $s = 300\,\mathrm{m}$? \hfill (2 BE)

\vspace{2cm}

\item[d)] Zeichne das $s(t)$-Diagramm für $t = 0$ bis $t = 10\,\mathrm{s}$. \hfill (2 BE)

\begin{center}
\begin{tikzpicture}
\begin{axis}[
    width=10cm, height=5cm,
    xlabel={$t$ in s},
    ylabel={$s$ in m},
    xmin=0, xmax=10,
    ymin=0, ymax=350,
    grid=major,
    axis lines=left,
    xtick={0,2,4,6,8,10},
    ytick={0,50,100,150,200,250,300,350}
]
\end{axis}
\end{tikzpicture}
\end{center}

\item[e)] Ein zweiter Zug startet zur gleichen Zeit am Bahnhofseingang ($s_0 = 0$) und beschleunigt mit $a = 2\,\mathrm{m/s^2}$. Wann überholt er den ersten Zug? \hfill (2 BE)

\vspace{2.5cm}
\end{enumerate}

\newpage

\aufgabe{Freier Fall mit Anfangsgeschwindigkeit}{10}

Ein Ball wird von einem 80 m hohen Turm mit $v_0 = 10\,\mathrm{m/s}$ senkrecht nach unten geworfen.

\begin{enumerate}
\item[a)] Berechne die Fallzeit bis zum Aufprall. \hfill (3 BE)

\vspace{3.5cm}

\item[b)] Berechne die Aufprallgeschwindigkeit. \hfill (2 BE)

\vspace{2.5cm}

\item[c)] Vergleiche mit dem freien Fall (ohne Anfangsgeschwindigkeit): Wie viel Zeit spart man durch das Werfen? \hfill (3 BE)

\vspace{3cm}

\item[d)] Erkläre, warum die Zeiteinsparung relativ gering ist, obwohl man den Ball mit 10 m/s wirft. \hfill (2 BE)

\vspace{2.5cm}
\end{enumerate}

\vspace{0.5cm}
\hrule
\vspace{0.3cm}

\noindent\fbox{\parbox{\dimexpr\textwidth-2\fboxsep-2\fboxrule}{%
\textbf{Zusatzaufgabe} (für schnelle Schüler, max. 2 Zusatz-BE)
}}

\vspace{0.3cm}

Zwei Straßenbahnen fahren auf parallelen Gleisen. Bahn A fährt mit konstant $v = 15\,\mathrm{m/s}$. Bahn B startet 100 m hinter Bahn A aus dem Stand und beschleunigt mit $a = 2\,\mathrm{m/s^2}$.

Nach welcher Zeit hat Bahn B die Bahn A eingeholt? \hfill (2 BE)

\vspace{3cm}

\hrule
\vspace{0.3cm}
\begin{center}
\textit{Erlaubte Hilfsmittel: Tafelwerk}\\[0.2cm]
\textit{Hinweis: $g \approx 10\,\mathrm{m/s^2}$}
\end{center}

\end{document}
